%% ma: L12-B15-0018
%% lop: 12
%% chuong: 5
%% bai: 15
%% dang: 12.15.14
%% loai: TN4
%% muc: VD
%% dap_an: "B"
%% nguon: "(NBV)-10 ĐỀ KIỂM TRA HK2-TOÁN 12-DUNG CHUNG 3 SÁCH.pdf (H02, Đề số 2), Phần 1 Câu 3"
%% kien_thuc: "Bài 15 (hình chiếu vuông góc của đường thẳng lên mặt phẳng), Bài 14"
%% trang_thai: cho-duyet
%% ly_do: "Dạng mới 12.15.14 chờ duyệt."
%% ghi_chu: "Đáp án gốc B khớp giải độc lập."
\begin{ex}
Trong không gian $Oxyz$, cho đường thẳng $d\colon\dfrac{x-1}{2}=\dfrac{y+5}{-1}=\dfrac{z-3}{4}$. Phương trình nào dưới đây là phương trình hình chiếu vuông góc của $d$ trên mặt phẳng $(P)\colon x+3=0$?
\choice{$\begin{cases}x=-3\\y=-5+2t\\z=3-t\end{cases}$}{\True $\begin{cases}x=-3\\y=-6-t\\z=7+4t\end{cases}$}{$\begin{cases}x=-3\\y=-5-t\\z=-3+4t\end{cases}$}{$\begin{cases}x=-3\\y=-5+t\\z=3+4t\end{cases}$}
\loigiai{Phương trình tham số của $d$: $x=1+2t$, $y=-5-t$, $z=3+4t$. Mặt phẳng $(P)\colon x=-3$ vuông góc với $Ox$, nên phép chiếu vuông góc lên $(P)$ giữ nguyên $y$, $z$ và thay $x$ bởi $-3$. Hình chiếu của $d$ là đường thẳng $\Delta\colon x=-3$, $y=-5-t$, $z=3+4t$, có vectơ chỉ phương $(0;-1;4)$.

Phương án B có vectơ chỉ phương $(0;-1;4)$ và đi qua $(-3;-6;7)$, ứng với $t=1$ của $\Delta$ nên trùng $\Delta$. Phương án A, D có vectơ chỉ phương $(0;2;-1)$, $(0;1;4)$ không cùng phương với $(0;-1;4)$. Phương án C đi qua $(-3;-5;-3)$ không thuộc $\Delta$. Chọn B.}
\end{ex}
